% Auto-generated from meeting_2569-01-30_monthly-jan.md (กบ.วช. format v2)
\documentclass{meetingminutes}
\meetingcommittee{คณะกรรมการบริหารหลักสูตร วิศวกรรมคอมพิวเตอร์และปัญญาประดิษฐ์ (บัณฑิตศึกษา)}
\meetingnumber{๓/ภาคเรียนที่ ๒/๒๕๖๘}{๒๕๖๘}
\meetingdate{วันที่ ๓๐ มกราคม พ.ศ. ๒๕๖๙}
\meetingplace{ห้องประชุมคณะเทคโนโลยีอุตสาหกรรม}
\meetingstart{๑๓.๐๐ น.}
\meetingend{๑๕.๑๕ น.}

\begin{document}
\maketitle

\psection{ผู้มาประชุม}
\begin{participants}
  \pmember{1}{ผู้ช่วยศาสตราจารย์ ดร.พิศิษฐ์ นาคใจ}{อาจารย์ประจำหลักสูตร}{ประธาน}
  \pmember{2}{ผู้ช่วยศาสตราจารย์ ดร.กาญจนา ดาวเด่น}{อาจารย์ประจำหลักสูตร (เลขานุการ)}{กรรมการ/เลขานุการ}
  \pmember{3}{ผู้ช่วยศาสตราจารย์ ดร.พัชรี มณีรัตน์}{อาจารย์ประจำหลักสูตร}{กรรมการ}
\end{participants}

\psection{ผู้ไม่มาประชุม}
\noindent ไม่มี\par

\startmeeting
\openingchair{ผู้ช่วยศาสตราจารย์ ดร.พิศิษฐ์ นาคใจ}{กล่าวเปิดการประชุมและดำเนินการประชุมตามระเบียบวาระ ดังนี้}

\agenda{๑}{รับรองรายงานการประชุมครั้งที่ ๒/ภาคเรียนที่ ๒/๒๕๖๘}
ที่ประชุมรับรองรายงานการประชุมครั้งที่ ๒/ภาคเรียนที่ ๒/๒๕๖๘ โดยไม่มีการแก้ไข

\noindent\textbf{มติที่ประชุม}\quad รับรองรายงานการประชุม

\agenda{๑.๑}{รายงานผลการประชุมคณะกรรมการบริหารหลักสูตรฯ ครั้งที่ ๑/๒๕๖๙ (๗ มกราคม ๒๕๖๙)}
ประธานแจ้งให้ที่ประชุมทราบว่า ที่ประชุมคณะกรรมการบริหารหลักสูตรฯ ครั้งที่ ๑/๒๕๖๙
(วันที่ ๗ มกราคม ๒๕๖๙) มีมติเห็นชอบการปรับปรุงแก้ไขหลักสูตรตาม สมอ. ๐๘
โดยขอเพิ่มอาจารย์ประจำหลักสูตรจำนวน ๓ คน ได้แก่
ผศ.ดร.พิทักษ์ คล้ายชม, ผศ.ดร.รัฐพล ดุลยะลา และ ผศ.ดร.วีระพล คงนุ่น
เพื่อเสริมความครอบคลุมด้านวิศวกรรมไฟฟ้า ระบบควบคุมอัตโนมัติ และการจัดการทรัพยากร
ขณะนี้อยู่ระหว่างการนำเสนอต่อคณะกรรมการบัณฑิตศึกษาและคณะกรรมการบริหารวิชาการประจำคณะต่อไป

\noindent\textbf{มติที่ประชุม}\quad ที่ประชุมรับทราบ


\agenda{๒}{รายงานผลการเรียนกลางภาค และ Academic Performance}
อาจารย์ผู้สอนรายงานผลการเรียนกลางภาค:

วิชา 4095101 — ปัญญาประดิษฐ์และการเรียนรู้ของเครื่อง: - สอบกลางภาคเสร็จสิ้น นักศึกษาทุกคนสอบผ่าน - CLO1-CLO3 (ทฤษฎี AI, การวิเคราะห์ประสิทธิภาพโมเดล) อยู่ในระดับดี - ไม่มีนักศึกษาที่มีความเสี่ยงด้านผลการเรียน

วิชา 4095102 — การเขียนโปรแกรมขั้นสูงสำหรับการเรียนรู้ของเครื่อง: - นักศึกษาทุกคนส่งโครงงานพัฒนาโมเดล ML ระยะแรกแล้ว คุณภาพโดยรวมดี - ผศ.ดร.พิศิษฐ์ รายงานว่า CLO4 (จริยธรรมในการพัฒนาโปรแกรม) ซึ่งสอดคล้องกับ PLO4 มีพัฒนาการที่ดี


\agenda{๓}{ติดตามความก้าวหน้าหัวข้อวิจัย/สารนิพนธ์}
อาจารย์ที่ปรึกษารายงานความก้าวหน้า: - นายเกริกเกียรติ — อยู่ระหว่างเก็บข้อมูล Dataset สำหรับ Deep Learning - นางสาวพอใจ — ออกแบบต้นแบบระบบ IoT เบื้องต้นแล้ว - นายจาตุรงค์ — ศึกษาวรรณกรรมด้าน AI ในระบบพลังงาน - นายจิรกิตติ์ — เก็บข้อมูล Sentinel-2 และเริ่มเทรนโมเดล ความก้าวหน้าดีที่สุดในกลุ่ม


\agenda{๔}{รายงานความก้าวหน้าอาจารย์ด้านการวิจัยและการพัฒนาวิชาชีพ}
ผศ.ดร.พิศิษฐ์ รายงานว่า: - ได้รับทุนวิจัย Fundamental Fund จาก สป.อว. สำหรับโครงการวิจัย AI - นายจิรกิตติ์ วราภรณ์ ได้รับรางวัลในการแข่งขันด้าน Remote Sensing ระดับภาคเหนือ (ชนะเลิศประเภทนำเสนอผลงาน และรางวัลทองแดงประเภทผลงานวิจัย) ซึ่งเป็นผลงานที่น่าภาคภูมิใจในฐานะนักศึกษารุ่นแรกของหลักสูตร


\adjournmeeting

\begin{sigblock}
\typist{ผู้ช่วยศาสตราจารย์ ดร.กาญจนา ดาวเด่น}
\reviewer{ผู้ช่วยศาสตราจารย์ ดร.พิศิษฐ์ นาคใจ}{กรรมการ}
\chairsig{ผู้ช่วยศาสตราจารย์ ดร.พิศิษฐ์ นาคใจ}{ประธานการประชุม}
\end{sigblock}

\end{document}
